\chapter{Methodology}
\label{sec:methodology}

This chapter describes how the proposed system works and how it will be evaluated.
The presentation is deliberately ordered from intuition to formalism: each mechanism
is first described in plain terms, then stated precisely. Every equation is followed
by a reading of what it says in words, so that the formal notation documents the
design rather than substituting for an explanation of it.

\section{The Central Idea}

A conventional Retrieval-Augmented Generation system behaves like a student who
visits the library exactly once. Whatever material is gathered in that single trip
is all the material available when the essay is written. If a claim turns out to
need a supporting table that was not collected, the student cannot go back --- and
the most likely outcome is that the missing number is invented rather than the gap
admitted.

The proposed system behaves instead like a careful researcher. It gathers some
evidence, stops to ask whether that evidence is actually enough, returns for more
when it is not, and refuses to write until the question can be answered from what it
has in hand. Three mechanisms make this behaviour concrete:

\begin{enumerate}
    \item An \textbf{agent} that decides, at each step, which kind of evidence to
    fetch next --- and may fetch repeatedly rather than once.
    \item An \textbf{Evidence Sufficiency Gate} that stands between gathering and
    writing, and only opens when the gathered evidence can support an answer.
    \item A \textbf{Faithfulness Gate} that checks the written answer against the
    gathered evidence, and sends it back for rewriting if any claim is unsupported.
\end{enumerate}

The remainder of this chapter formalises these mechanisms, describes the
retrieval tools available to the agent, explains how evidence from different
modalities is combined, and defines the evaluation plan.

\section{Agentic System Architecture}

\subsection{The System as a State Machine}

The agent is built on the ReAct paradigm~\cite{yao2023react}, in which a model alternates between
reasoning about what it needs and acting to obtain it. Formally, the system is a
state machine:
\begin{equation}
    \mathcal{M} = \langle \mathcal{S}, \mathcal{A}, \mathcal{T}, \mathcal{O}, \mathcal{E} \rangle
\end{equation}

\noindent In words, the system is fully described by five things: what it currently
knows, what it is able to do, how doing something changes what it knows, what it
gets back when it acts, and how it decides when to stop. Each component is defined
below.

\begin{description}
    \item[$\mathcal{S}$ --- the state.] Everything the system currently knows. This
    comprises the interpreted intent of the query, the record of which retrievals
    have already been performed, and the evidence buffer $\mathcal{B}_t$ holding
    every piece of evidence gathered so far.

    \item[$\mathcal{A}$ --- the action space.] The set of tools the agent may
    invoke. Each tool retrieves a different modality.

    \item[$\mathcal{T}$ --- the transition function.] How the state changes once a
    tool has been used: new evidence is added to the buffer and the retrieval
    history is extended.

    \item[$\mathcal{O}$ --- the observation space.] What comes back from a tool ---
    a text passage, a parsed table, or a page region.

    \item[$\mathcal{E}$ --- the Evidence Sufficiency Gate.] The decision procedure
    that determines whether gathering should continue or generation may begin.
\end{description}

\subsection{Query Decomposition}

Given a complex scientific query $Q$, the agent does not attempt to satisfy it in
one retrieval. It first breaks the query into a sequence of smaller, individually
answerable sub-goals:
\begin{equation}
    Q \longrightarrow \{q_1, q_2, \dots, q_k\}
\end{equation}

\noindent For example, the query \emph{``Does the accuracy gain reported in
Section~4 hold for the low-resource subset?''} decomposes into two sub-goals: first
locate the claim about the accuracy gain in the narrative text, then locate the
table that reports results broken down by subset. Decomposition matters because the
second sub-goal cannot even be formulated until the first has been answered --- it is
the retrieved claim that tells the agent which table to look for. This is what is
meant by \emph{multi-hop} reasoning, and it is the reason a single retrieval pass is
structurally incapable of answering such questions.

\subsection{The Reasoning Loop}

At each time step $t$ the agent selects one action $a_t$ from the action space:
\begin{equation}
    \mathcal{A} = \{\texttt{TextRetriever}, \texttt{TableRetriever}, \texttt{FigureRetriever}, \texttt{RegionZoom}\}
\end{equation}

\noindent Executing $a_t$ returns an observation $o_t$, which is appended to the
evidence buffer. After $t$ steps the buffer therefore contains the accumulated
result of every retrieval performed so far:
\begin{equation}
    \mathcal{B}_t = \bigcup_{i=1}^{t} o_i
\end{equation}

\noindent The union notation simply says that nothing is discarded: evidence
retrieved at step 1 is still available at step 5. The agent then consults the
Sufficiency Gate, and either loops back to select another tool or proceeds to
generation.

\subsection{The Evidence Sufficiency Gate}

The Sufficiency Gate is the principal novelty of the architecture, and its purpose
is best stated as a question the system asks itself before writing anything:
\emph{``Can I answer this from what I have actually retrieved?''}

Concretely, a meta-evaluator inspects the query $Q$ alongside the current buffer
$\mathcal{B}_t$ and returns a judgement of sufficiency. Consider the running
example. After the first retrieval the buffer contains the narrative sentence
asserting an accuracy gain, but nothing about the low-resource subset. The claim is
present; the evidence qualifying it is not. The gate therefore fails, the agent
reformulates and dispatches the \texttt{TableRetriever}, and the gate is
re-evaluated against the enlarged buffer.

The significance of this design is that it is \emph{preventative}. A conventional
pipeline would forward the incomplete buffer to the generator and depend on the
generator to notice what was missing --- which is precisely the situation that
produces fluent, confident, unsupported answers. The gate removes the opportunity to
make that error rather than attempting to detect it afterwards.

The loop is bounded by a maximum iteration count. If sufficiency is not reached
before the bound is exhausted, the system reports that the evidence is inadequate
rather than answering anyway. For a literature review tool this is the correct
failure behaviour: a researcher can act on a stated gap, but cannot detect a
fabrication.

\subsection{The Faithfulness Gate}

The Sufficiency Gate guarantees that the evidence needed for an answer is present;
it does not guarantee that the generator uses only that evidence. A second gate
therefore runs after generation. The draft answer is split into individual factual
claims, and each claim is checked against the evidence buffer $\mathcal{B}_t$. If
any claim is not supported by an item in the buffer, the gate fails and the
generator is asked to rewrite the unsupported part using only buffered evidence.

Because the generator has no access to the corpus other than through
$\mathcal{B}_t$, this check is bounded: each claim is compared against a fixed,
finite set of retrieved items rather than against the world at large. Like the
retrieval loop, the rewrite loop has a maximum iteration count; if it is exhausted,
unsupported claims are removed and reported as unconfirmed rather than kept.

\section{Modality-Specific Retrieval Tools}

Scientific PDFs are not uniform documents, and a single retrieval mechanism handles
them poorly. The action space therefore provides four tools, three matched to a
different kind of content and one for inspecting a page in finer detail. Each is described below together with the reason it
exists as a separate tool.

\begin{description}
    \item[\texttt{TextRetriever}.] Performs dense semantic retrieval over chunked
    narrative text using SPECTER2~\cite{singh2023scirepeval}, a model built on
    SciBERT~\cite{beltagy2019scibert} and trained specifically for scientific
    document retrieval. A retrieval-trained model is required because a
    pre-trained language model such as SciBERT is not, on its own, trained to
    place a query close to the passage that answers it; and a scientific model is
    used rather than a general-purpose one because scientific vocabulary is
    precisely where general embeddings are weakest.

    \item[\texttt{TableRetriever}.] Works in two steps. First, it locates candidate
    tables by querying the visual (ColPali) index and keeping pages whose
    highest-scoring region is a table. Second, it runs a table-structure recogniser,
    such as the Table Transformer~\cite{smock2022pubtables}, on that region; this
    recovers rows, columns and a bounding box for every cell, and each cell's text
    is read from the PDF's text layer at that position. The result is a normalised
    table \emph{that preserves the two-dimensional coordinates of every cell}. Coordinate preservation is not incidental detail. The bounding boxes
    required by Objective~4 can only be produced if the spatial layout of the table
    survives parsing; a parser that returned a table as flat text would satisfy the
    reasoning requirement and silently destroy the attribution requirement.

    \item[\texttt{FigureRetriever} (ColPali).] Embeds entire PDF pages as grids of
    visual patches using ColPali~\cite{faysse2025colpali}, a vision-language model
    that scores pages by late interaction between query tokens and page
    patches~\cite{khattab2020colbert}, locating charts,
    plots and diagrams without depending on OCR. This matters because OCR fails
    most severely exactly where scientific documents are most information-dense ---
    multi-column headers, merged cells, and axis labels. Treating the page as an
    image sidesteps that failure mode entirely.

    \item[\texttt{RegionZoom}.] Crops a sub-region of a page already returned by
    the \texttt{FigureRetriever} and re-encodes it at higher resolution. A
    page-level match identifies which page holds a table or chart, but not which
    row or data point answers the question; zooming supplies that finer level of
    detail.
\end{description}

\section{Cross-Modal Evidence Fusion}

Once evidence has been retrieved from different modalities, the system must relate
the pieces to one another. A symbol appearing in a sentence and the same symbol
appearing in a chart legend must be recognised as referring to the same quantity;
otherwise the system holds two disconnected facts rather than one linked piece of
evidence.

This linkage is performed in two places, neither of which requires training a new
model component.

\begin{enumerate}
    \item \textbf{Explicit linking in the evidence buffer.} Every item added to
    $\mathcal{B}_t$ carries its provenance: source paper, page, retrieving tool,
    and step. Items are additionally linked when they share a page or when one
    explicitly references the other --- for example, a sentence containing
    ``Table~3'' is linked to the table whose caption begins ``Table~3''. The buffer
    is therefore a small graph of connected evidence rather than an unordered list.

    \item \textbf{Joint reasoning in the generator.} The linked items are passed to
    the vision-language generator as a single interleaved context, with each text
    passage placed next to the table or figure crops it is linked to. The
    generator's own attention layers then relate text tokens to image regions ---
    this is the joint text--image reasoning such models are pre-trained to perform,
    and no additional fusion layer is added on top of it.
\end{enumerate}

The practical consequence is that downstream generation is conditioned on linked
multimodal evidence rather than on two independent evidence streams that happen to
have been retrieved together.

\section{Model Choice}

Token attribution by integrated gradients (Section~\ref{sec:explainability})
requires computing gradients of the generator's output with respect to its input,
which is only possible with full access to the model's weights. The generator is
therefore an open-weight vision-language model run locally, such as
Qwen2-VL~\cite{wang2024qwen2vl}, rather than a model reached through a hosted API.
For simplicity, the same model also serves as the planner and as the evaluator
behind both gates.

\section{Grounded Multimodal Explainability}
\label{sec:explainability}

A system that produces a correct answer without showing where it came from remains
unusable for academic work, because a reader has no way to check it. The generation
phase therefore emits three audit artefacts alongside the natural-language response.
Each answers a different verification question.

\begin{enumerate}
    \item \textbf{Text Token Attribution} --- \emph{``Which sentences did this come
    from?''} Integrated gradients~\cite{sundararajan2017axiomatic}, computed on the open-weight generator, identify the source sentences that most influenced
    each generated claim, which are returned as highlights in the original text.

    \item \textbf{Visual Saliency (Bounding Boxes)} --- \emph{``Where on the page is
    the proof?''} Explicit coordinates $[x_{\min}, y_{\min}, x_{\max}, y_{\max}]$
    mark the exact table cell or chart region the claim rests on. For tables, these
    come directly from the cell coordinates preserved by the \texttt{TableRetriever};
    for figures, they are obtained by thresholding the ColPali patch-similarity map
    computed by the \texttt{FigureRetriever} (or \texttt{RegionZoom}) and taking
    the enclosing box of the highest-scoring patches. This is the artefact that
    discharges Objective~4.

    \item \textbf{Retrieval DAG Trace} --- \emph{``How did the system get here?''} A
    logged execution graph recording which tools were called, in what order, and how
    many times each gate rejected the buffer or the draft before opening. This makes
    the reasoning path auditable rather than merely asserted.
\end{enumerate}

Together these mean a reader need not trust the system's answer. They can check it
against the highlighted sentence, look at the boxed table cell, and inspect the path
by which both were found.

\section{Experimental Setup and Evaluation Plan}

\subsection{Dataset}

The framework will be evaluated on \textbf{SPIQA}~\cite{pramanick2024spiqa}, a dataset of over 270{,}000
questions over computer-science papers whose answers depend on the papers' figures
and tables. SPIQA is appropriate here because each question is tied to a reference
figure or table, so the system must locate specific visual evidence rather than
answer from narrative text alone. Each question is annotated with that reference
figure or table, which provides figure-level ground truth for retrieval.

SPIQA does not annotate \emph{regions} within a figure or table. Evaluating
bounding-box precision (CGS, below) therefore requires a manually annotated subset
in which the supporting cell or chart region is marked for each question.
Similarly, since not every SPIQA question is multi-hop, results will also be
reported separately on the subset of questions that require combining text with a
figure or table, which is where the difference between iterative and single-pass
retrieval is expected to be largest.

\subsection{Baselines}

Two baselines isolate the two contributions of this work:

\begin{itemize}
    \item \textbf{Single-Pass Multimodal RAG} (e.g.\ a LLaVA-based~\cite{liu2023llava}
    pipeline that retrieves once, then answers) --- multimodal but
    not iterative. Comparison against this baseline isolates the contribution of
    \emph{agency}: what is gained by being permitted to retrieve more than once.
    \item \textbf{Text-Only Agentic Reasoning} (e.g.\ ReAct~\cite{yao2023react}) --- iterative but not
    multimodal. Comparison against this baseline isolates the contribution of
    \emph{visual retrieval}: what is gained by being able to see tables and figures
    at all.
\end{itemize}

\noindent Because each baseline holds one contribution fixed and removes the other,
an improvement over both cannot be attributed to either factor alone.

\subsection{Evaluation Metrics}

Four metrics are proposed. Accuracy alone is insufficient here, because a system may
reach a correct answer through unverifiable means --- and such an answer does not
satisfy the aims of this work.

\begin{description}
    \item[Retrieval Faithfulness (RF).] The precision of the retrieved evidence
    against the ground-truth evidence for each question. A low score indicates the
    system reached its answer without ever holding the material that justifies it.
    \emph{Addresses RQ1 and RQ2.}

    \item[Modality Attribution Precision (MAP).] How accurately generated content is
    traced back to the correct source modality. A low score indicates the system
    cannot reliably distinguish what it read in the text from what it read in a
    figure. \emph{Addresses RQ2.}

    \item[Agent Efficiency Score (AES).] The ratio of necessary retrieval hops to
    total tool calls. This guards against a degenerate solution in which the agent
    achieves high accuracy by retrieving indiscriminately; such a system would be
    accurate but neither efficient nor interpretable. \emph{Addresses RQ1.}

    \item[Cross-modal Grounding Score (CGS).] Intersection-over-Union between
    generated bounding boxes and ground-truth regions, measuring how precisely the
    system localises its visual evidence. Computed on the manually annotated subset
    described above. \emph{Addresses RQ3.}
\end{description}

\noindent Taken together, RF and AES measure whether the agentic loop works, MAP
measures whether cross-modal retrieval works, and CGS measures whether the
attribution is precise enough to be genuinely verifiable by a human reader.